\documentclass[12pt]{article}
\usepackage{amsmath, amssymb, amsthm}
\usepackage[margin=1in]{geometry}

\newtheorem{theorem}{Theorem}
\newtheorem{lemma}[theorem]{Lemma}

\title{Problem 261: Pivotal Square Sums}
\author{Project Euler Solution}
\date{}

\begin{document}
\maketitle

\section{Problem Statement}

A positive integer~$k$ is called a \textbf{square-pivot} if there exist integers $m>0$ and $n\ge k$ such that
\[
(k-m)^2 + (k-m+1)^2 + \cdots + k^2 = (n+1)^2 + (n+2)^2 + \cdots + (n+m)^2.
\]
Find the sum of all distinct square-pivots $k \le 10^{10}$.

\section{Mathematical Development}

Denote $S(t) = \sum_{i=1}^{t} i^2 = \tfrac{t(t+1)(2t+1)}{6}$.

\begin{theorem}[Reduction to generalized Pell equation]
\label{thm:pell}
If $k$ is a square-pivot with parameters $(m,n)$, then writing $m = sp^2$ ($s$ squarefree) and $m+1 = gr^2$ ($g$ squarefree), there exist non-negative integers $q,u$ satisfying
\[
gu^2 - sq^2 = 1.
\]
Equivalently, with $D = sg$, $X = gu$, $Y = q$: $X^2 - DY^2 = g$.
\end{theorem}

\begin{proof}
The pivot condition $S(k) - S(k-m-1) = S(n+m) - S(n)$ expands using the closed form of~$S$. Introducing $x = 2k - m$ and $t = 2n + m + 1$, algebraic simplification yields
\[
\frac{x^2}{m} - \frac{t^2}{m+1} = 1,
\]
which factors as $x^2 = m\cdot a$ and $t^2 = (m+1)(a+1)$ for some integer $a \ge 0$. With the squarefree decompositions $m = sp^2$ and $m+1 = gr^2$, we obtain $x = spq$ (giving $a = sq^2$) and $t = gru$ (giving $a+1 = gu^2$). Subtracting: $gu^2 - sq^2 = 1$. Setting $X = gu$, $Y = q$ gives
\[
X^2 - DY^2 = g^2u^2 - sg\,q^2 = g(gu^2 - sq^2) = g. \qedhere
\]
\end{proof}

\begin{lemma}[Base solution]
\label{lem:base}
The pair $(X_0, Y_0) = (gr, p)$ satisfies $X_0^2 - DY_0^2 = g$.
\end{lemma}

\begin{proof}
$X_0^2 - DY_0^2 = g^2r^2 - sg\,p^2 = g(gr^2 - sp^2) = g\bigl((m+1) - m\bigr) = g$.
\end{proof}

\begin{lemma}[Pell recurrence]
\label{lem:recurrence}
Let $(x_1, y_1)$ be the fundamental solution to $X^2 - DY^2 = 1$. Then all solutions to $X^2 - DY^2 = g$ in the orbit of $(X_0, Y_0)$ are generated by
\[
(X_{j+1}, Y_{j+1}) = (X_j\,x_1 + D\,Y_j\,y_1,\; X_j\,y_1 + Y_j\,x_1).
\]
\end{lemma}

\begin{proof}
By the Brahmagupta--Fibonacci identity:
\[
(X_jx_1 + DY_jy_1)^2 - D(X_jy_1 + Y_jx_1)^2 = (X_j^2 - DY_j^2)(x_1^2 - Dy_1^2) = g \cdot 1 = g.
\]
The new solution is strictly larger since $x_1 \ge 1$ and $y_1 \ge 1$, and the orbit exhausts all positive solutions in its class (a standard result in the theory of Pell equations).
\end{proof}

\begin{lemma}[Recovery of $k$ and $n$]
Given a Pell solution $(X, Y)$ with $Y = q$:
\[
k = \frac{sp(p+q)}{2}, \qquad n = \frac{gru - m - 1}{2} \quad\text{where } u = X/g.
\]
The pivot is valid if and only if $k \in \mathbb{Z}_{>0}$ and $n \ge k$.
\end{lemma}

\begin{proof}
From $x = spq$ and $x = 2k - m$, we get $k = (spq + sp^2)/2 = sp(p+q)/2$. From $t = gru$ and $t = 2n+m+1$, we get $n = (gru-m-1)/2$.
\end{proof}

\begin{lemma}[Bound on $m$]
For a valid pivot $k \le 10^{10}$, we have $m \le 70710$.
\end{lemma}

\begin{proof}
The constraint $n \ge k$ forces $k \ge 2m(m+1)$. Then $2m(m+1) \le 10^{10}$ gives $m^2 < 5\times 10^9$, so $m \le 70710$.
\end{proof}

\section{Algorithm}

\begin{verbatim}
pivots = empty set
for m = 1 to 70710:
    decompose m = s*p^2 (s squarefree), m+1 = g*r^2 (g squarefree)
    D = s*g; skip if D is a perfect square
    (x1, y1) = fundamental Pell solution for X^2 - D*Y^2 = 1
    (X, Y) = (g*r, p)   [base solution]
    loop:
        k = s*p*(p + Y) / 2
        if k > 10^10: break
        if k is a positive integer and k >= 2*m*(m+1):
            n = (g*r*(X/g) - m - 1) / 2
            if n >= k: pivots.add(k)
        (X, Y) = (X*x1 + D*Y*y1, X*y1 + Y*x1)
return sum(pivots)
\end{verbatim}

\section{Complexity Analysis}

\begin{itemize}
\item \textbf{Time:} $O(M\log\log M + M(\sqrt{D_{\max}} + \log N))$ where $M = 70710$ is the bound on~$m$, the $M\log\log M$ term accounts for the sieve, $\sqrt{D_{\max}}$ bounds the continued-fraction period for each Pell solve, and $\log N$ bounds the Pell iterates per orbit (solutions grow exponentially).
\item \textbf{Space:} $O(M + |\mathrm{pivots}|)$ for the sieve, Pell cache, and result set.
\end{itemize}

\section{Answer}

\[
\boxed{238890850232021}
\]

\end{document}
